\clearpage
\section[Introduction and Objectives]{Chapter -- 01: Introduction and Objectives}

This project reproduces a profiled, machine-learning side-channel attack on the ASCAD
dataset and quantifies how well four model families recover one byte of a masked
AES-128 key, using the number of attack traces needed to drive the true key byte to
rank zero as the primary measure of success.

\subsection{Background}

A masking countermeasure splits every sensitive intermediate value into random shares
so that no single point in a power trace depends on the secret key alone. ASCAD
(ANSSI, 2018)~\cite{benadjila2020} is the reference benchmark for evaluating whether a
profiled, learning-based attacker can still recover the key from such an
implementation. Poudel and Rahimi~\cite{poudel2025} frame the attack as supervised
256-class classification of the first-round AES S-box output for one key byte and
report key recovery within a few tens of attack traces on ASCAD. This project treats
that paper's pipeline as a target to rebuild end to end, and asks: which of its
qualitative claims still hold when every stage -- data loading, preprocessing,
training, and evaluation -- is written from scratch and checked against the real
ASCAD files.

\subsection{Objectives}

\begin{enumerate}
\item To build the full ASCAD key-recovery pipeline from the raw HDF5 files through
      to a Key Rank / Guessing Entropy evaluation, rather than relying on a
      pre-processed benchmark harness.
\item To reproduce the paper's central qualitative claim: that deep models (CNN,
      ResNet) recover the masked key byte while classical baselines (random forest,
      SVM) do not.
\item To implement and compare four model families -- random forest, SVM, a compact
      1-D CNN, and a 1-D ResNet -- under one shared training and evaluation
      interface.
\item To evaluate every model on both the fixed-key and variable-key ASCAD sets, and
      confirm the paper's finding that a ResNet is more trace-efficient than a CNN on
      the harder variable-key set.
\item To go beyond the paper with three extensions: a multi-task ResNet that also
      predicts the mask shares, a desynchronisation robustness test, and a
      shuffled-label sanity control.
\item To document every deviation from the paper and every engineering fix required
      to make the pipeline train reliably, so the whole benchmark is reproducible
      from the repository alone.
\end{enumerate}

\subsection{Data and Model}

Table~\ref{tab:data-model} summarises the target, the datasets, and the evaluation
metric used throughout the project. This defines the exact object every later chapter
operates on: a fixed target key byte, attacked under a masking countermeasure, scored
by how many traces it takes to rank that byte first.

\begin{table}[h]
\centering
\caption{Data and model used throughout the project.}
\label{tab:data-model}
\begin{tabular}{@{}p{3.2cm}p{10.5cm}@{}}
\toprule
Item & Description \\
\midrule
Target & Key byte 2 of an AES-128 key, via $y=\mathrm{Sbox}(p[2]\oplus k[2])$
         (256-class identity leakage model). ASCAD only masks bytes 2--15, and byte 2
         is the community-standard target. \\[4pt]
Datasets & Fixed-key ASCAD (ATMega8515): 50\,000 profiling + 10\,000 attack traces of
           700 samples. Variable-key ASCAD-r: 200\,000 + 100\,000 traces of 1400
           samples. Desynchronised variants shift each trace randomly by up to 50 or
           100 samples. \\[4pt]
Models & Random forest (300 trees), SVM (RBF kernel, PCA-50), a compact 1-D CNN
          (Zaid et al.)~\cite{zaid2020}, a 1-D ResNet (Karayalcin et
          al.)~\cite{karayalcin2023}, and a multi-task variant of the ResNet. \\[4pt]
Metric & Key Rank / Guessing Entropy: the rank of the true key byte among all 256
         hypotheses, as a function of the number of attack traces used, averaged over
         100 random trace orderings. Rank 0 means the key byte is recovered. \\
\bottomrule
\end{tabular}
\end{table}

As the table shows, top-1 classification accuracy is not the metric of interest here
-- on masked ASCAD it sits near the $1/256\approx0.39\%$ chance line for every model.
Key rank is what every subsequent section reports.

\textbf{Result:} The project's target key byte, its two datasets, and the Key Rank
evaluation metric were fixed at the outset and used consistently throughout every
later section of this report.
